\documentclass[a4paper,12pt]{article}
\usepackage[utf8]{inputenc}
\usepackage[T1]{fontenc}
\usepackage[ngerman]{babel}
\usepackage{amsmath,amssymb}
\usepackage{booktabs}
\usepackage{enumitem}
\usepackage[margin=2.5cm]{geometry}
\setlength{\emergencystretch}{2em}

\begin{document}

\begin{center}
{\LARGE\bfseries \"Ubungsblatt 5 -- L\"osungen}\\[0.5em]
{\large VC Logik}
\end{center}

\bigskip

%% ============================================================
%% AUFGABE 1
%% ============================================================
\subsection*{Aufgabe 1 \textnormal{(2 Punkte)} -- DPLL f\"ur eine CNF}

\textbf{Was ist eine CNF?}
CNF hei\ss t \emph{konjunktive Normalform}. Das bedeutet: Die ganze Formel
ist ein gro\ss es UND von mehreren Klammern. Jede Klammer ist eine ODER-Verkn\"upfung.
So eine einzelne Klammer nennt man \textbf{Klausel}. Ein einzelner Baustein
in einer Klausel hei\ss t \textbf{Literal}, also z.\,B. $X$ oder $\neg X$.

\[
(X \vee Y) \wedge (\neg X \vee Y)
\]
ist also in CNF, weil zwei Klauseln mit $\wedge$ verbunden sind:
\[
\{X,Y\}
\qquad\text{und}\qquad
\{\neg X,Y\}.
\]
Damit eine CNF wahr ist, muss jede Klausel wahr sein. Und eine Klausel ist
schon wahr, sobald mindestens ein Literal darin wahr ist.

\medskip
\textbf{Was ist DPLL?}
DPLL ist ein Verfahren, mit dem gepr\"uft wird, ob eine CNF erf\"ullbar ist.
``Erf\"ullbar'' bedeutet: Es gibt mindestens eine Belegung der Variablen,
bei der alle Klauseln gleichzeitig wahr werden.

Die Idee ist einfach: DPLL setzt Schritt f\"ur Schritt Werte f\"ur Variablen
ein und vereinfacht danach die Klauseln.
\begin{itemize}[nosep]
\item Wird eine Klausel wahr, ist sie erledigt und kann gestrichen werden.
\item Wird ein Literal falsch, wird nur dieses Literal aus seiner Klausel entfernt.
\item Bleibt in einer Klausel nur ein Literal \"ubrig, muss dieses Literal wahr werden.
      Das nennt man \textbf{Unit Propagation}.
\item Wird eine Klausel ganz leer, ist dieser Zweig unm\"oglich. Dann muss
      DPLL zur\"uckgehen und einen anderen Wert probieren.
\end{itemize}
Wenn keine Unit-Klausel vorhanden ist, wird eine Variable ausgew\"ahlt und
ein Wert probiert. Genau das passiert hier gleich mit $X$.

\medskip
Gegeben ist
\[
F \equiv (X \vee Y) \wedge (X \vee \neg Y)
\wedge (\neg X \vee Y) \wedge (\neg X \vee \neg Y).
\]
Die Formel ist schon in konjunktiver Normalform. Die vier Klauseln
werden so nummeriert. Rechts steht jeweils noch die Mengen-Schreibweise
f\"ur DPLL; der Beistrich in der Menge bedeutet hier also ODER innerhalb
der Klausel:
\[
\begin{array}{rclcl}
C_1 &\equiv& (X \vee Y)        &=& \{X,Y\},\\
C_2 &\equiv& (X \vee \neg Y)   &=& \{X,\neg Y\},\\
C_3 &\equiv& (\neg X \vee Y)   &=& \{\neg X,Y\},\\
C_4 &\equiv& (\neg X \vee \neg Y) &=& \{\neg X,\neg Y\}.
\end{array}
\]

\medskip
\textbf{Start.}
Es gibt am Anfang keine Unit-Klausel und auch kein reines Literal.
Also muss DPLL eine Variable verzweigen. Als erste Variable wird $X$
gew\"ahlt.

\begin{center}
\small
\renewcommand{\arraystretch}{1.25}
\begin{tabular}{p{0.17\textwidth}p{0.22\textwidth}p{0.51\textwidth}}
\toprule
Schritt & Belegung & Was mit den Klauseln passiert \\
\midrule
Start
& noch keine Werte
& $C_1,C_2,C_3,C_4$ bleiben alle offen. \\
\midrule
Wahl $X=1$
& $X=1$
& $C_1=(X\vee Y)$ wird wahr, weil $X=1$ ist. Also wird $C_1$
  gestrichen.\\
&& $C_2=(X\vee\neg Y)$ wird aus demselben Grund wahr. Also wird auch
  $C_2$ gestrichen.\\
&& In $C_3=(\neg X\vee Y)$ ist $\neg X$ falsch, weil $X=1$ ist.
  Der falsche Teil $\neg X$ wird entfernt; es bleibt nur $\{Y\}$.\\
&& In $C_4=(\neg X\vee\neg Y)$ wird ebenfalls $\neg X$ entfernt;
  es bleibt $\{\neg Y\}$.\\
&& Rest nach diesem Schritt: $\{Y\}$ und $\{\neg Y\}$. \\
\midrule
Unit Prop.
& $X=1,\;Y=1$
& $\{Y\}$ ist eine Unit-Klausel. Damit diese Klausel wahr wird, muss
  $Y=1$ gesetzt werden.\\
&& Dann ist aber $\neg Y$ falsch. In der Klausel $\{\neg Y\}$ wird
  also das einzige Literal entfernt. Dadurch bleibt die leere Klausel
  $\emptyset$.\\
&& Eine leere Klausel bedeutet: Diese Klausel kann nicht mehr wahr
  gemacht werden. Der Zweig $X=1$ scheitert. \\
\bottomrule
\end{tabular}
\end{center}

Der Zweig $X=1$ funktioniert also nicht. DPLL geht zur\"uck und probiert
den anderen Wert:

\begin{center}
\small
\renewcommand{\arraystretch}{1.25}
\begin{tabular}{p{0.17\textwidth}p{0.22\textwidth}p{0.51\textwidth}}
\toprule
Schritt & Belegung & Was mit den Klauseln passiert \\
\midrule
Wahl $X=0$
& $X=0$
& Jetzt ist $\neg X=1$. Deshalb werden $C_3=(\neg X\vee Y)$ und
  $C_4=(\neg X\vee\neg Y)$ wahr und gestrichen.\\
&& In $C_1=(X\vee Y)$ ist $X$ falsch, also wird $X$ entfernt;
  es bleibt $\{Y\}$.\\
&& In $C_2=(X\vee\neg Y)$ wird ebenfalls $X$ entfernt;
  es bleibt $\{\neg Y\}$.\\
&& Rest nach diesem Schritt: wieder $\{Y\}$ und $\{\neg Y\}$. \\
\midrule
Unit Prop.
& $X=0,\;Y=1$
& Aus $\{Y\}$ folgt wieder $Y=1$.\\
&& Dann ist $\neg Y$ wieder falsch. Die Klausel $\{\neg Y\}$ verliert
  ihr einziges Literal und wird zu $\emptyset$.\\
&& Also scheitert auch der Zweig $X=0$. \\
\bottomrule
\end{tabular}
\end{center}

\medskip
\noindent\fbox{\parbox{\dimexpr\textwidth-2\fboxsep-2\fboxrule}{%
\textbf{Ergebnis Aufgabe~1.}\\[2pt]
Beide m\"oglichen Zweige f\"ur $X$ enden in einem Widerspruch.
Also ist
\[
(X \vee Y) \wedge (X \vee \neg Y)
\wedge (\neg X \vee Y) \wedge (\neg X \vee \neg Y)
\]
\textbf{unerf\"ullbar}. Es gibt keine Belegung f\"ur $X,Y$, die alle
vier Klauseln gleichzeitig wahr macht.
}}

\newpage

%% ============================================================
%% AUFGABE 2
%% ============================================================
\subsection*{Aufgabe 2 \textnormal{(2 Punkte)} -- Erst CNF, dann DPLL}

Gegeben ist
\[
F \equiv \neg X \vee \neg Y \to \neg(\neg Y \vee X).
\]
Zuerst muss klar sein, wo die Klammern sitzen. Die Implikation $\to$
bindet schw\"acher als $\vee$. Daher wird die Formel so gelesen:
\[
F \equiv (\neg X \vee \neg Y) \to \neg(\neg Y \vee X).
\]
Links von $\to$ steht also der Block $(\neg X \vee \neg Y)$,
rechts von $\to$ steht der Block $\neg(\neg Y \vee X)$.

\medskip
\textbf{Schritt 1: In CNF bringen.}
F\"ur DPLL braucht man eine CNF. Deshalb wird die Formel zuerst so
umgeformt, dass sie aus Klauseln besteht. Die erste Regel ist:
\[
P \to Q \equiv \neg P \vee Q.
\]
Hier ist
\[
P=(\neg X \vee \neg Y)
\qquad\text{und}\qquad
Q=\neg(\neg Y \vee X).
\]
Also wird aus $P \to Q$:
\begin{align*}
F
&\equiv \neg(\neg X \vee \neg Y) \vee \neg(\neg Y \vee X).
\end{align*}

\smallskip
Jetzt stehen vor zwei Klammern Negationen. Diese werden mit De Morgan
nach innen gezogen. Die Regel lautet:
\[
\neg(A \vee B) \equiv \neg A \wedge \neg B.
\]
F\"ur den ersten Teil bedeutet das:
\begin{align*}
\neg(\neg X \vee \neg Y)
&\equiv \neg(\neg X) \wedge \neg(\neg Y)\\
&\equiv X \wedge Y.
\end{align*}
F\"ur den zweiten Teil:
\begin{align*}
\neg(\neg Y \vee X)
&\equiv \neg(\neg Y) \wedge \neg X\\
&\equiv Y \wedge \neg X.
\end{align*}
Damit wird die ganze Formel:
\begin{align*}
F
&\equiv (X \wedge Y) \vee (Y \wedge \neg X).
\end{align*}

\smallskip
Jetzt sieht man in beiden Teilen ein $Y$:
\[
(X \wedge Y)
\qquad\text{und}\qquad
(Y \wedge \neg X).
\]
Also kann $Y$ ausgeklammert werden. Das ist dieselbe Idee wie bei
Zahlen: $3a + 3b = 3(a+b)$. Nur wird hier mit $\wedge$ und $\vee$
gerechnet:
\begin{align*}
(X \wedge Y) \vee (Y \wedge \neg X)
&\equiv Y \wedge (X \vee \neg X).
\end{align*}
Der Teil $(X \vee \neg X)$ ist immer wahr: Entweder ist $X$ wahr oder
$\neg X$ ist wahr. Deshalb bleibt im Grunde nur $Y$:
\[
F \equiv Y.
\]
Das hei\ss t: Die ganze urspr\"ungliche Formel verh\"alt sich genau wie
die einfache Formel $Y$. Die CNF besteht deshalb nur aus einer Klausel:
\[
\mathcal{K} = \bigl\{\{Y\}\bigr\}.
\]

\medskip
\textbf{Schritt 2: DPLL anwenden.}
Jetzt ist der DPLL-Teil sehr kurz. Es gibt nur noch die Klausel $\{Y\}$.
Das ist eine Unit-Klausel, weil nur ein Literal darin steht.

\begin{center}
\small
\renewcommand{\arraystretch}{1.25}
\begin{tabular}{p{0.18\textwidth}p{0.22\textwidth}p{0.50\textwidth}}
\toprule
Schritt & Belegung & Was passiert \\
\midrule
Start
& noch keine Werte
& Die einzige Klausel ist $\{Y\}$. Damit diese Klausel wahr wird,
  muss $Y$ wahr sein. \\
\midrule
Unit Prop.
& $Y=1$
& $\{Y\}$ ist erf\"ullt und wird gestrichen. Danach bleibt keine
  offene Klausel mehr \"ubrig. DPLL hat also ein Modell gefunden. \\
\bottomrule
\end{tabular}
\end{center}

F\"ur $X$ muss DPLL nichts mehr entscheiden. Der Grund ist: Nach der
Umformung steht $X$ gar nicht mehr in der Formel. Also ist $X$ frei.
Solange $Y=1$ ist, funktioniert sowohl $X=0$ als auch $X=1$.

\medskip
\noindent\fbox{\parbox{\dimexpr\textwidth-2\fboxsep-2\fboxrule}{%
\textbf{Ergebnis Aufgabe~2.}\\[2pt]
Die Formel ist \"aquivalent zu $Y$ und damit \textbf{erf\"ullbar}.
Eine passende Belegung ist zum Beispiel
\[
Y=1,\qquad X=0.
\]
Genauso w\"urde auch $Y=1,\;X=1$ funktionieren, weil $X$ am Ende keine
Rolle mehr spielt.
}}

\newpage

%% ============================================================
%% AUFGABE 3
%% ============================================================
\subsection*{Aufgabe 3 \textnormal{(2 Punkte)} -- Logische Konsequenz mit DPLL}

Gegeben ist
\begin{align*}
\mathrm{KB}_1 \equiv {}&
(A \vee B \vee \neg C)
\wedge (\neg A \vee B)
\wedge (A \vee C \vee \neg B) \\
&\wedge (D \vee \neg C)
\wedge \neg A.
\end{align*}
Gefragt ist, ob
\[
\mathrm{KB}_1 \models (C \wedge B)
\]
gilt.

\medskip
\textbf{Schritt 1: Die Frage f\"ur DPLL umformen.}
Eine logische Konsequenz pr\"uft man \"uber ein Gegenbeispiel:
\[
\mathrm{KB}_1 \models G
\quad\Longleftrightarrow\quad
\mathrm{KB}_1 \wedge \neg G \text{ ist unerf\"ullbar}.
\]
Hier ist $G = C \wedge B$. Also:
\[
\neg(C \wedge B) \equiv \neg C \vee \neg B.
\]
Daher wird mit DPLL getestet, ob folgende CNF erf\"ullbar ist:
\[
\Phi \equiv \mathrm{KB}_1 \wedge (\neg C \vee \neg B).
\]
Als Klauselmenge:
\[
\begin{array}{rclcrcl}
C_1 &=& \{A,B,\neg C\},       && C_2 &=& \{\neg A,B\},\\
C_3 &=& \{A,C,\neg B\},       && C_4 &=& \{D,\neg C\},\\
C_5 &=& \{\neg A\},           && C_6 &=& \{\neg C,\neg B\}.
\end{array}
\]

\medskip
\textbf{Schritt 2: DPLL starten.}
Es gibt sofort die Unit-Klausel $C_5=\{\neg A\}$.
Also muss $A=0$ gelten.

\begin{center}
\renewcommand{\arraystretch}{1.25}
\begin{tabular}{p{0.2\textwidth}p{0.3\textwidth}p{0.38\textwidth}}
\toprule
Schritt & Belegung & Rest nach Vereinfachung \\
\midrule
Unit $C_5$
& $A=0$
& $C_2$ und $C_5$ sind wahr.
  Aus $C_1$ wird $\{B,\neg C\}$.
  Aus $C_3$ wird $\{C,\neg B\}$.
  $C_4=\{D,\neg C\}$ und $C_6=\{\neg C,\neg B\}$ bleiben. \\
\bottomrule
\end{tabular}
\end{center}

Nach $A=0$ bleibt also:
\[
\{\{B,\neg C\},\;\{C,\neg B\},\;\{D,\neg C\},\;\{\neg C,\neg B\}\}.
\]
Jetzt gibt es keine Unit-Klausel. Als n\"achste Verzweigung wird $C$
gew\"ahlt.

\medskip
\textbf{Zweig 1: $C=1$.}

\begin{center}
\renewcommand{\arraystretch}{1.25}
\begin{tabular}{p{0.2\textwidth}p{0.3\textwidth}p{0.38\textwidth}}
\toprule
Schritt & Belegung & Rest nach Vereinfachung \\
\midrule
Wahl $C=1$
& $A=0,\;C=1$
& $\{C,\neg B\}$ ist wahr.
  Aus $\{B,\neg C\}$ wird $\{B\}$.
  Aus $\{\neg C,\neg B\}$ wird $\{\neg B\}$.
  Aus $\{D,\neg C\}$ wird $\{D\}$. \\
\midrule
Unit Propagation
& $A=0,\;C=1,\;B=1$
& $\{B\}$ erzwingt $B=1$.
  Dann wird aber $\{\neg B\}$ falsch, also entsteht
  $\emptyset$. Dieser Zweig scheitert. \\
\bottomrule
\end{tabular}
\end{center}

Der Wert $C=1$ passt also nicht. DPLL geht zur\"uck und probiert $C=0$.

\medskip
\textbf{Zweig 2: $C=0$.}

\begin{center}
\renewcommand{\arraystretch}{1.25}
\begin{tabular}{p{0.2\textwidth}p{0.3\textwidth}p{0.38\textwidth}}
\toprule
Schritt & Belegung & Rest nach Vereinfachung \\
\midrule
Wahl $C=0$
& $A=0,\;C=0$
& $\{B,\neg C\}$, $\{D,\neg C\}$ und $\{\neg C,\neg B\}$
  sind wahr, weil $\neg C$ wahr ist.
  Aus $\{C,\neg B\}$ bleibt $\{\neg B\}$. \\
\midrule
Unit Propagation
& $A=0,\;C=0,\;B=0$
& $\{\neg B\}$ erzwingt $B=0$.
  Danach sind alle Klauseln erf\"ullt. F\"ur $D$ ist kein Wert mehr
  erzwungen. \\
\bottomrule
\end{tabular}
\end{center}

Damit liegt ein Modell f\"ur
$\Phi = \mathrm{KB}_1 \wedge \neg(C\wedge B)$ vor:
\[
A=0,\qquad B=0,\qquad C=0,\qquad D=0
\]
(auch $D=1$ w\"are m\"oglich).

\medskip
\textbf{Kurzer Check des Gegenbeispiels.}
Mit $A=0,B=0,C=0,D=0$ gilt:
\[
\begin{array}{rcl}
A \vee B \vee \neg C &=& 0 \vee 0 \vee 1 = 1,\\
\neg A \vee B &=& 1 \vee 0 = 1,\\
A \vee C \vee \neg B &=& 0 \vee 0 \vee 1 = 1,\\
D \vee \neg C &=& 0 \vee 1 = 1,\\
\neg A &=& 1.
\end{array}
\]
Also ist $\mathrm{KB}_1$ wahr. Aber
\[
C \wedge B = 0 \wedge 0 = 0.
\]
Das ist genau ein Gegenbeispiel gegen die behauptete Folgerung.

\medskip
\noindent\fbox{\parbox{\dimexpr\textwidth-2\fboxsep-2\fboxrule}{%
\textbf{Ergebnis Aufgabe~3.}\\[2pt]
\[
\mathrm{KB}_1 \not\models (C \wedge B).
\]
Der Grund ist: DPLL findet f\"ur
$\mathrm{KB}_1 \wedge \neg(C \wedge B)$ ein Modell, n\"amlich z.\,B.\
\[
A=0,\quad B=0,\quad C=0,\quad D=0.
\]
Damit ist die Wissensbasis wahr, aber $C \wedge B$ falsch.
}}

\end{document}
